\subsection*{Câu 1}

\noindent\textbf{1. Mạch kiến thức:} Quy tắc dấu ngoặc

\noindent\textbf{2. Độ khó:} *

\noindent\textbf{3. Đề bài:}
\[
-\frac{11}{24}+\left(\frac{5}{41}-\frac{13}{24}\right)
\]

\noindent\textbf{4. Lời giải:}
\begin{align*}
-\frac{11}{24}+\left(\frac{5}{41}-\frac{13}{24}\right)
    &= -\frac{11}{24}+\frac{5}{41}-\frac{13}{24} \\
    &= \frac{5}{41}-\left(\frac{13}{24}-\frac{11}{24}\right) \\
    &= \frac{5}{41}-\frac{2}{24} \\
    &= \frac{5}{41}-\frac{1}{12}=\frac{13}{84}
\end{align*}

\noindent\textbf{5. Token bắt đầu sai:}
\[
-\frac{11}{24}
\]

\noindent\textbf{Dòng bắt đầu sai:} 3

\noindent\textbf{6. Giải thích lỗi:}

Khi thay đổi vị trí các số hạng trong biểu thức, phải giữ nguyên dấu của từng số hạng.

\noindent\textbf{7. Sửa lại:}

\textit{Dạng câu hỏi:} Trắc nghiệm

\textit{Câu hỏi:} Chọn phép biến đổi đúng sau khi sửa lỗi.


 \begin{itemize}
    \item[A.] $\displaystyle \frac{5}{41}+\left(\frac{11}{24}-\frac{13}{24}\right)$
    \item[B.] $\displaystyle \frac{5}{41}-\left(\frac{13}{24}+\frac{11}{24}\right)$
    \item[C.] $\displaystyle \frac{5}{41}+\left(\frac{13}{24}-\frac{11}{24}\right)$
    \item[D.] $\displaystyle \frac{5}{41}-\left(\frac{13}{24}-\frac{11}{24}\right)$
\end{itemize}
\textit{Đáp án:} B

\subsection*{Câu 2}

\noindent\textbf{1. Mạch kiến thức:} Quy tắc dấu ngoặc

\noindent\textbf{2. Độ khó:} **

\noindent\textbf{3. Đề bài:}
\[
\left(-\frac{3}{4}+\frac{5}{13}\right)\cdot\frac{7}{2}
-\left(\frac{9}{4}-\frac{8}{13}\right)\cdot\frac{7}{2}
\]

\noindent\textbf{4. Lời giải:}
\begin{align*}
\left(-\frac{3}{4}+\frac{5}{13}\right)\cdot\frac{7}{2}
-\left(\frac{9}{4}-\frac{8}{13}\right)\cdot\frac{7}{2}
&=
\left(-\frac{3}{4}+\frac{5}{13}-\frac{9}{4}-\frac{8}{13}\right)\cdot\frac{7}{2} \\
&=
\left(-\frac{3}{4}-\frac{9}{4}+\frac{5}{13}-\frac{8}{13}\right)\cdot\frac{7}{2} \\
&=
\left(-3-\frac{3}{13}\right)\cdot\frac{7}{2} \\
&=
-\frac{42}{13}\cdot\frac{7}{2} \\
&=
-\frac{147}{13}
\end{align*}

\noindent\textbf{5. Token bắt đầu sai:}
\[
-\frac{8}{13}
\]

\noindent\textbf{6. Giải thích lỗi:}

Khi bỏ dấu ngoặc có dấu ``$-$'' đằng trước, phải đổi dấu tất cả các số hạng trong dấu ngoặc.

\noindent\textbf{7. Sửa lại:}

\textit{Dạng câu hỏi:} Trắc nghiệm

\noindent\textbf{7. Sửa lại:}

\textit{Dạng câu hỏi:} Trắc nghiệm

\textit{Câu hỏi:}
\[
\left(-\frac{3}{4}+\frac{5}{13}-\frac{9}{4}+\frac{8}{13}\right)\cdot\frac{7}{2}
\]

Chọn cách gộp các số hạng hợp lí để tiếp tục tính.

\begin{itemize}
    \item[A.] $\displaystyle
    \left[\left(-\frac34-\frac94\right)
    +\left(\frac5{13}+\frac8{13}\right)\right]\cdot\frac72$

    \item[B.] $\displaystyle
    \left[\left(-\frac34+\frac94\right)
    +\left(\frac5{13}+\frac8{13}\right)\right]\cdot\frac72$

    \item[C.] $\displaystyle
    \left[\left(-\frac34+\frac5{13}\right)
    +\left(\frac94+\frac8{13}\right)\right]\cdot\frac72$

    \item[D.] $\displaystyle
    \left[\left(-\frac34-\frac94\right)
    -\left(\frac5{13}+\frac8{13}\right)\right]\cdot\frac72$
\end{itemize}

\textit{Đáp án:} A

\subsection*{Câu 3}

\noindent\textbf{1. Mạch kiến thức:} Quy tắc dấu ngoặc

\noindent\textbf{2. Độ khó:} ***

\noindent\textbf{3. Đề bài:}
\[
-\frac{5}{7}\cdot\frac{2}{11}
-\frac{5}{7}\cdot\frac{9}{11}
+\frac{5}{7}
\]

\noindent\textbf{4. Lời giải:}
\begin{align*}
-\frac{5}{7}\cdot\frac{2}{11}
-\frac{5}{7}\cdot\frac{9}{11}
+\frac{5}{7}
&=
-\frac{5}{7}\left(\frac{2}{11}+\frac{9}{11}+1\right) \\
&=
-\frac{5}{7}\cdot\frac{20}{11} \\
&=
-\frac{100}{77}
\end{align*}

\noindent\textbf{5. Token bắt đầu sai:}
\[
+1
\]

\noindent\textbf{6. Giải thích lỗi:}

Khi đặt $-\frac{5}{7}$ làm thừa số chung, dấu của mỗi số hạng trong ngoặc phải được xác định sao cho khi nhân phân phối với $-\frac{5}{7}$, ta thu lại đúng biểu thức ban đầu.\\
\noindent\textbf{7. Sửa lại:}

\textit{Dạng câu hỏi:} Trắc nghiệm

\textit{Câu hỏi:} Sau khi sửa lỗi, ta có:
\[
-\frac{5}{7}\left(\frac{2}{11}+\frac{9}{11}-1\right)
\]

Chọn kết quả của biểu thức bằng cách tính toán hợp lí.

\begin{itemize}
    \item[A.] $-\frac{5}{7}$
    \item[B.] $0$
    \item[C.] $\frac{5}{7}$
    \item[D.] $-\frac{10}{7}$
\end{itemize}

\textit{Đáp án:} B

\subsection*{Câu 4}

\noindent\textbf{1. Mạch kiến thức:} Quy tắc chuyển vế

\noindent\textbf{2. Độ khó:} **

\noindent\textbf{3. Đề bài:}
\[\frac{3}{10}x-1\frac12 =\left(-\frac27\right):\frac5{14}\]

\noindent\textbf{4. Lời giải:}
\begin{align*}
\frac{3}{10}x-1\frac{1}{2}&=\left(-\frac{2}{7}\right):\frac{5}{14}\\
\frac{3}{10}x-\frac{3}{2}&=-\frac{4}{5}\\
\frac{3}{10}x&=-\frac{4}{5}+\frac{3}{2}\\
\frac{3}{10}x&=\frac{7}{10}\\
x&=\frac{7}{10}-\frac{3}{10}\\
x&=\frac{2}{5}
\end{align*}

\noindent\textbf{5. Token bắt đầu sai:}
\[
-
\]

\noindent\textbf{Dòng bắt đầu sai:} 5

\noindent\textbf{6. Giải thích lỗi:}

Muốn tìm một thừa số chưa biết trong một tích, ta lấy tích chia cho thừa số đã biết.

\noindent\textbf{7. Sửa lại:}

\textit{Dạng câu hỏi:} Trắc nghiệm

\textit{Câu hỏi:} Sau khi sửa lỗi, phép tính tìm $x$ ở bước cuối là:

\begin{itemize}
    \item[A.] $\displaystyle x=\frac7{10}-\frac3{10}$
    \item[B.] $\displaystyle x=\frac7{10}+\frac3{10}$
    \item[C.] $\displaystyle x=\frac7{10}:\frac3{10}$
    \item[D.] $\displaystyle x=\frac7{10}\cdot\frac3{10}$
\end{itemize}

\textit{Đáp án:} C

\subsection*{Câu 5}

\noindent\textbf{1. Mạch kiến thức:} Thứ tự thực hiện các phép tính

\noindent\textbf{2. Độ khó:} **

\noindent\textbf{3. Đề bài:}
\[
B=4,75+\left(-\frac12\right)^3+0,5^2-3\cdot\left(-\frac38\right)
\]

\noindent\textbf{4. Lời giải:}
\begin{align*}
B
&=4,75+\left(-\frac12\right)^3+0,5^2-3\cdot\left(-\frac38\right)\\
&=4,75+\frac18+0,25+\frac98\\
&=(4,75+0,25)+\left(\frac18+\frac98\right)\\
&=5+\frac{10}{8}\\
&=5+\frac54\\
&=6,25
\end{align*}
\noindent\textbf{5. Token bắt đầu sai:}
\[
\frac18
\]

\noindent\textbf{6. Giải thích lỗi:}

Khi nâng một số âm lên lũy thừa với số mũ lẻ, kết quả vẫn là số âm.

\noindent\textbf{7. Sửa lại:}

\textit{Dạng câu hỏi:} Trắc nghiệm

\textit{Câu hỏi:} Sau khi sửa lỗi, chọn phép biến đổi đúng để tiếp tục tính:
\[
B=4,75-\frac18+0,25+\frac98
\]

\begin{itemize}
    \item[A.] $\displaystyle B=(4,75+0,25)+\left(-\frac18+\frac98\right)$
    \item[B.] $\displaystyle B=(4,75-0,25)+\left(-\frac18+\frac98\right)$
    \item[C.] $\displaystyle B=(4,75+0,25)+\left(\frac18+\frac98\right)$
    \item[D.] $\displaystyle B=(4,75+0,25)-\left(-\frac18+\frac98\right)$
\end{itemize}

\textit{Đáp án:} A